% Version 3.9: original DATA progress plus explicitly scoped request inclusion.
% This fragment does not repeat the prefix-safety or program-isolation proofs.
\subsection{Progress premises and stable retries}
\label{sec:v35-progress-premises}
The results below use the frozen model and the specified retry,
publication, retention and public-kernel rules. They do not assume that
the next required DATA is accepted during a successful service interval
in order to prove that acceptance. Publication still starts with the
original complete regional result. \Cref{prop:carrier-preservation}
extracts its existing public kernel from the logical carrier, including
a prior bad-hint record; it does not promise a new retry or weaker
publication input. Initial authority, historical-byte hydration and
native state preparation remain eventual, with no uniform delay bound.
The quantitative source argument concerns prepared regional processing
that witnesses canonical $G$ progress, not every observer's wall clock.

\paragraph{Inherited regional conditions.}
The v3.1 progress paragraph states: ``The regional engine keeps its
original fair-schedule, exclusion and uninterrupted-commit-window
conditions.'' The retained Part~I conditions~\cite{part1} therefore
include
\[
 \forall g_0\ \exists g\ge g_0:\ L(g,0)\text{ is correct},\qquad
 \forall g,t_0\ \exists t\ge t_0:\ L(g,t)\text{ is correct}.
\]
For every exposed equivocator $p$, a finite generation $G_p$ excludes
$p$ from all leader positions at generations $g\ge G_p$. These are
generation indices; exclusion changes future authorities, not an
already finalized assignment. The additional condition in
\cref{sec:progress}, concerning correct turns that are actually entered,
is retained. A name in a schedule is not evidence that its turn has
entered or its DATA has become global.

For one installed slot $r$, write its immutable schedule as
\[
 L(\beta,r,g,t)=\ell_{r,g}(i_{\beta,r},t).
\]
Changing $\beta$ does not reseed this schedule. A zero-output retry reads
the same $\tagg{AUTH}_{r,g}$ records. It creates fresh locks, votes,
entry times and timers for its new pair and initially enters $g=t=0$.
Only a Quit proof from this new pair admits its later recoveries.
In particular, an old-beta Quit is not a new-beta entry proof.

Suppose, for a contradiction, that $G$ eventually stays at $G^*$.
At most one subsequent Force can close a duty that had already added
positive DATA to this frontier. After it, the zero-output rule keeps a
fixed slot $r^*$. Its configuration, quota and base-window parameters
remain fixed. Persistent preparation retains the executable state of
this exact anchor, its provenance and validated dependencies; persistent
authority fetching retains both cached records and unfinished requests
across beta changes. The node-level store preserves raw namespaces,
including rejected evidence, and completes the fixed-cut hydration of
\cref{sec:source-hydration}. The background authority task continues through
successive indices of $r^*$, rather than restarting at zero on Force.

All finite sets used below belong to this installed configuration and
the existing finite region cycle. No finite union of the memberships
of every future slot, and no uniform delay bound over those future
configurations, is assumed.

\subsection{A finite regional response after preparation}
\label{sec:v35-regional-response}
\begin{proposition}[Finite authority prefix and regional response]
\label{prop:v35-regional-response}
Fix $r^*$ and a region $i$ under its installed snapshot. In a tail with
fixed $G^*$, there is a finite recovery index $g_i^*$ and a finite
preparation time after which all authorities through $g_i^*$ and the
exact opening state are available from the persistent cache. Consider
a fresh retry after that time, with no external closure during the
interval under consideration. Under the complete ReadyBundle handler,
it forms and offers a complete result containing positive session DATA
within a finite $R_i$ of its first correct Enter.
\end{proposition}
\begin{proof}
The installed regional roster is finite. Let $E_i$ be the set of its
identities whose admissible equivocation is ever exposed in this fixed
authority namespace. For each $p\in E_i$, its first exposure gives a
finite exclusion generation $G_{i,p}$. Since all beta retries use the
same immutable schedules, that exclusion tail applies to each retry.
Set
\[
 \bar g_i=\max\bigl(\{1\}\cup\{G_{i,p}:p\in E_i\}\bigr),\qquad
 g_i^*=\min\{g\ge\bar g_i:
                   \ell_{r^*,g}(i,0)\text{ is correct}\}.
\]
The finite roster makes the maximum finite; inherited initializer
fairness makes $g_i^*$ finite. No admissible equivocation can be exposed
in a generation at or above $\bar g_i$: its signer would belong to
$E_i$ while occupying a leader position forbidden by its exclusion
tail. A proof naming a lower generation or another pair does not freeze
this generation. Low-generation faulty assignments can recur on a new
beta, but the number of generations before $g_i^*$ is fixed.

The required set $\{\tagg{AUTH}_{r^*,g}:0\le g\le g_i^*\}$ is finite.
Eventual authority availability and the persistent fetching task make
all its records available after a finite time. The same applies to the
fixed anchor's native preparation and validation dependencies. The fixed
historical cut contains finitely many source objects; eventual retrieval
and continued retention hydrate it after a finite time. New namespaces
are retained while that fixed task runs, rather than extending its
startup target indefinitely. Taking a maximum of these finite completion
times does not impose a uniform archive, authority or state-restoration
bound. Cached executable
anchor state is attached directly to a fresh continuation; an unbounded
database import on every retry would not satisfy this preparation rule.

Here is an explicit conservative regional bound after preparation.
Put $Q=\TB_{r^*}$, $\Delta=\Delta_i$, and
\[
 D_i=(Q+1)(2K+12)\Delta_i,\qquad
 R_i=(g_i^*+1)D_i+2g_i^*\Delta_i+9\Delta_i.
\]
Initial correct entries differ by at most $\Delta$. The same bound
holds for a recovery with cached authority; ordinary term-entry skew
is at most $2\Delta$, by Part~I's terminal-QC handoff argument. The
REPORT timing argument in \cref{prop:v38-report-timing} preserves the original
first-DATA voting deadlines: at most $4\Delta$ for an ordinary term and
$6\Delta$ for term zero. Reports create no new progress clock, and the
correct initializer's ANCHOR vote remains immediate. These are entry
events of the current pair, not inherited timestamps.

Let $u$ be the first correct entry to a term. Its last progress check
occurs by $u+(2K+6)\Delta$, unless a Quit, Completion or advancement
has already intervened. If the opening is still unknown, the first
check runs anyway; missing that opening cannot suppress BLAME forever.
If a correct member has voted for the terminal proposal by this bound,
complete forwarding and correct votes yield its QC within another
$2\Delta$, unless a Quit intervenes. Below the endpoint this advances
the term. At the endpoint, an uninterrupted interval of another
$4\Delta$ covers all correct commits and collection of their FINAL
messages. If no correct member has cast that terminal vote, no correct
member can have advanced using its QC. The correct members' last
checks then produce matching BLAME; within another $2\Delta$ they all
freeze. Equivocation can cause an earlier Quit, not an unbounded wait
in this argument. Thus $(2K+12)\Delta$ bounds a term's advancement,
freeze or complete-result outcome.

A DATA generation has at most $Q$ terms, and a DRAIN generation has
one. Therefore $D_i$ bounds the time until a generation either forms a
complete result or reaches a correct Quit. After a first correct Quit,
its proof propagates within $\Delta$, and all correct members complete
the original recovery wait within $2\Delta$ of that first Quit when
the next authority is cached. Applying this argument to the finite
prefix reaches $g_i^*$ within the allowance in $R_i$, unless a positive
complete result already exists. This proves actual entry to that
initializer; schedule fairness alone was not used as a substitute for
entry.

At $g_i^*$ the initializer is correct and no equivocation can interrupt
the generation. Part~I's protected initial commit window gives its
first fresh proposal by $e+5\Delta$, all correct commitments by
$e+8\Delta$, and a complete FINAL receipt by $e+9\Delta$, where $e$
is first entry to this generation. If the fresh block is DRAIN, its
receipt completes the session's terminal DATA instead. Since $Q\ge1$,
that receipt still covers positive DATA after the retry's fixed anchor.
This proves the stated bound. Existing Completion only makes the
complete-result outcome earlier.
\end{proof}

\subsection{From an opening block to global progress}
\label{sec:v35-global-progress}
\begin{proposition}[Response witnessed by prepared regional replay]
\label{prop:v35-public-response}
In a fixed-$G^*$ tail, after finite preparation and the last effective
health-control change, let a retry open in a block first finalized at
real time $a$. For the continuously online, hydrated members of the
installed regional snapshots, there is a finite
\[
 L=\max_{i\in\mathcal I^*}
       \{2T_E+2\Delta_i+R_i+T_{\rm pub}\}
\]
such that, by $a+L$, their specified processing has witnessed canonical
$G$ growth or an effective Force has already closed this retry. Here
$\mathcal I^*$ is the finite set of regions selectable from the installed
snapshot. This is not a bound on arbitrary nonregional observers'
application processing or initial catch-up.
\end{proposition}
\begin{proof}
Take a finite maximum of the preparation times in
\cref{prop:v35-regional-response}, including fixed-cut hydration. Keep
the verified replay checkpoint, raw public bytes and directly executable
anchor objects. By \cref{prop:retained-public-bytes}, the opening block
and every predecessor finalized no later than it have their required
Espresso bytes available by $a+T_E$. Previously acquired raw evidence
has not become an archive cache miss. Controls are available before
their crossings, and the fixed slot's initial authority and exact anchor
provenance are prepared.

We check every kind of operation before the first positive DATA:
\begin{enumerate}[label=(\roman*),itemsep=2pt]
\item A canonical prelude uses the complete applicable controls,
retained source timestamps and existing window/accounting state. With
zero output the retry target remains $r^*$, whose initial authority
is prepared. A pending selection uses that same fixed authority and
retained eligibility/owed snapshots. Neither operation asks for an
arbitrary candidate's recovery authority.
\item Structural inventory reads the actual logical objects and
required public DATA/FINAL references. Inline bytes come from the
complete new namespace; historical Espresso aliases resolve in the
retained raw dictionary, including rejected records. A genuinely absent
required object gives the same incomplete-kernel decision at every
observer. Optional L1/admission/native sidecar targets are not fetched
to decide that kernel or to advance source inventory.
\item Type, hash, current-pair, fixed-context and signature checks
reject malformed or unrelated roots using the prepared initial snapshot.
Independently represented neighboring kernels are still inspected. A
root passing \tagg{CheckFINAL} has intrinsic regional input validity by
\cref{prop:final-kernel-soundness}. The publisher's alternative PROPOSE
or $\tagg{AUTH}_{r^*,g'}$ path cannot require fetching $g'>g_i^*$.
The kernel does not recheck an arbitrary L1-reference or native sidecar;
those intrinsic predicates were checked in the correct signer's genuine
regional execution. Exact DATA bytes and the independent global gates
remain mandatory.
\item Current append permission uses the actual fixed boundary, exact
$G^*$ and supplied hash-linked path, quota and source position. All
required Espresso references are retained and checked for strict
precedence. No current-pair DATA transaction imports native state or
runs the separate $C_{\rm pub}$ execution/DA publication operation.
Those interfaces remain sound and separate.
\item Initial Enter attaches the already executable exact anchor and
uses the prepared full Ready inputs. Regional Vote, recovery and receipt
formation still use full validation bundles and the cached finite AUTH
prefix; \cref{prop:v35-regional-response} covers their own processing.
A terminal first DATA commits and exposes $G$ before any successor
$\tagg{AUTH}_{r^*+1,0}$ acquisition.
\end{enumerate}
Thus all mandatory source dependencies before first DATA are either
retained preparation, new bounded Espresso inputs, advance-announced
controls, or original complete regional messages. No optional proof
choice creates an uncharged external lookup. Actual missing source or
boundary evidence would still be a wait outside this prepared case;
it is never silently reclassified as empty history.

The immediate \tagg{SOURCESTEP} relay and complete handler therefore
give first Ready/Enter by $a+T_E+\Delta_i$, unless verified closure
already forbids it. The ReadyBundle entry skew is part of $R_i$ and is
not an additional uncounted stage. This uses the maintained input path
and the complete-message interface stated in
\cref{sec:source-step-timing}, not a constant-time catch-up theorem for
arbitrary historical backlogs. Ready requires the opening header, not
a subsequently generated block.

\Cref{prop:v35-regional-response} now supplies the original complete
regional result containing positive DATA of this retry. Let $s$ be its
first correct offer time, with
$s\le a+T_E+\Delta_i+R_i$, and let $B_1$ be its first DATA after $G^*$.
Its direct or covering FINAL root covers $B_1$ in this same session.
The complete result supplies the DATA/ancestry and evidence required by
the unchanged publication input. A later-session receipt is not used
as a substitute, and a weaker kernel-only submission is not promised.

Consider its first logical carrying record, whether submitted by a
correct worker or previously by an adversary. If it was not already
recorded at $s$, the frozen delivery bounds make that record available
by $s+T_{\rm pub}+T_E$. If it was already finalized by $s$, its source
input is available by $s+T_E$, or earlier. No corrected wrapper is
needed. By \cref{prop:carrier-preservation}, the record's actual
structured result preserves the public kernel despite wrong outer
hints, an ancestry-only location or unrelated bad evidence. This is a
logical-value statement, not a binary-codec compatibility assertion.

The carrying record follows the retry's authenticated opening. A
correct FINAL signer committed only after validating a finalized
Espresso reference at or after that opening. Its signature, and hence
the complete result, exists only afterward. Immutability then places
the carrying block strictly after its reference and, for a forced retry,
after its opening block. This causal order is independent of publisher
hints. A correct signer cannot supply a same-block self-reference before
that block finalizes.

If $G$ has already grown or a Force has closed the retry, the claimed
alternative holds. Otherwise its canonical position still has exact
$G^*$ and this retry's current duty. The source case split above applies
also to every preceding record. The carrying record's kernel supplies
$B_1$ at the next height and passes \tagg{CheckFINAL} and the independent
public gates. The complete regional source handler commits and exposes
$B_1$ within its final $\Delta_i$ allowance. If another valid candidate
committed first, growth already holds; regional safety excludes a
different committed DATA at that height of this authorized session.

For quota one, $B_1$ is terminal. Its DATA, counts and pending REGULAR
boundary commit before \tagg{ResolveBoundary} obtains the next authority
or runs \tagg{PrepareOwed}. The outer record need not finish before this
growth becomes visible. A Force in the carrying block itself executes
before payload and supplies the closure alternative. No completed old
record is revisited to rescue a closed version.

Adding the opening source and relay, regional response, original
publication service, and result source and relay gives the displayed
$L$. The two $\Delta_i$ terms are the stated complete source handlers;
no new archive, L1 AUTH, native-output or external-observer processing
bound is hidden in either term.
\end{proof}

\begin{proposition}[Continuing global DATA progress]
\label{prop:v35-global-progress}
Under the frozen regional, source and delivery premises and the specified
retry, retention, public-kernel, complete regional source-handler and
health-controller rules, the canonical global DATA height grows without bound whenever the continuing eligibility condition
leaves a nonempty set of installed regions able to serve under those
premises. The bound does not require two eligible regions. It supplies
no uniform failover time for initial restoration or authority waits.
\end{proposition}
\begin{proof}
Assume instead that $G$ eventually equals a fixed $G^*$. The stable-slot
rule fixes $r^*$ after at most one positive-duty transition. The
health-controller policy permits only finitely many further matching
MSIG interventions and health-status changes at this unchanged frontier;
already submitted controls still take effect under their original
target, priority and first-crossing rules. Discard that finite control
prefix and complete the finite preparation above.

With fixed governance and authority eligibility, early cooldowns expire
at finite positions because Espresso timestamps grow. The finite owed
set stabilizes as well; if a round has finished, its next boundary opens
the next set from the governing eligible regions. Nonempty continuing
eligibility supplies a candidate. The conditional AUTO cooldown does
not remove the last eligible candidate in the current owed set: it
cools the source only when another such candidate exists. If the source
is outside that set, an eligible owed candidate already exists. Thus,
after this finite prefix, each effective AUTO immediately assigns the
next duty. There is no AUTO-created unassigned interval that could
replenish arbitrary real-time clock lag. This argument neither clears
an existing cooldown nor overrides OFF.

Number these consecutive retry openings by $n=0,1,\ldots$, and write
their first finalization times and timestamps as $a_n,\theta_n$.
If a retry is never closed, \cref{prop:v35-public-response} makes $G$
grow within $L$, a contradiction. Otherwise each closing event is AUTO
and is the next opening event. Every attempt has zero accepted DATA,
so it remains in the startup phase. Let $M_n=1+\beta_n$ be its fixed
tolerance and $W_n$ its unchanged startup base-window length. The
finite installed parameter set gives $W_{\min}=\min_n W_n>0$.
Settling crossed base windows yields
\[
 \theta_{n+1}-\theta_n\ge M_nW_n\ge M_nW_{\min},\qquad
 a_{n+1}-a_n\le L.
\]
The second inequality follows from
\cref{prop:v35-public-response} under the assumed absence of growth;
immediate reassignment is essential to it. A timestamp jump may settle
several misses, but causes only one Force and opens the new duty at the
realized timestamp. Since each effective Force increments beta,
$M_n=1+\beta_0+n$. Hence
\[
 \theta_0+W_{\min}\sum_{n<N}(1+\beta_0+n)
 \le\theta_N
 \le a_N+\epsilon_E
 \le a_0+NL+\epsilon_E.
\]
The left side is quadratic in $N$, while the right side is linear.
This contradiction uses only the frozen upper-lead bound. Any finite
initial timestamp lag is part of the constant $a_0-\theta_0$; no
lower-lag or empty-block-rate assumption was added. Thus $G$ cannot
have a last accepted DATA height.
\end{proof}

The conclusion concerns regions satisfying the stated regional model,
with the required correct voting weight and communication service. A
physically healthy region kept permanently OFF is not an eligible
candidate. Nor does this result extend the model to permanently missing
correct quorums in other assigned snapshots. The finite maxima above
are proof witnesses for a stalled frontier, not operational delay
parameters shared by every future configuration.

\subsection{Useful service for persistently eligible regions}
\label{sec:v35-regional-service}
The public round ledger records an owed set $\mathcal O$. At a round
opening it includes the region IDs that are governance-ON with positive
designated-authority weight; temporary AUTO cooldown does not remove a
member from this set. A region leaves the current set when its DATA is
accepted into $G$, or when effective governance OFF or loss of authority
eligibility prunes it. An ON after the round opened joins the next
round. A new round begins only after the old debts are discharged or
pruned and a subsequent duty boundary or selection opens it. The
ledger cannot create infinitely many rounds within one finite duty.

Only Force selection is filtered by $\mathcal O$. REGULAR selection
retains its historical \tagg{NextLive} rule over all eligible regions.
Its new lower bound is
\[
 q_0^{\rm regular}=
 \max\{\operatorname{height}(\xi(T_{\beta,r})),
                         q^{\rm open}_{\beta,r}+1\}.
\]
Thus accepted consecutive regular sessions have strictly increasing
opening heights. A later OFF does not rewrite an already fixed regular
witness. The ledger restricts the next Force choice; it does not revoke
such witnesses or retroactively change an open session's quota.

\begin{proposition}[Repeated accepted service]
\label{prop:v35-regional-service}
Assume the conditions of \cref{prop:v35-global-progress}. A region $C$
that remains governance-ON with positive designated-authority weight
from some position onward, and continues to satisfy the regional
serving premises, contributes DATA to $G$ repeatedly. Temporary AUTO
cooldown does not discharge this obligation.
\end{proposition}
\begin{proof}
First fix a round in which $C$ is owed and suppose that it never
contributes. Its debt cannot be pruned. Every other member of this
finite round can be discharged or pruned only once; new ON events do
not add members to the unfinished round.

If effective Forces occur infinitely often, consider the first accepted
DATA in each Force-created version that produces any global DATA.
It must come from that version's initially selected, owed region: the
strict cutoff rejects prior-version output, and contiguous acceptance
cannot skip the forced session. This first DATA either discharges that
member or the member was already pruned since selection. With $C$
unserved, only finitely many such removals are possible. Infinitely many
productive Force-created versions are therefore impossible. On the
other hand, if only finitely many versions after those Forces produce
DATA, continuing effective Forces eventually leave a tail with no
global growth, contradicting \cref{prop:v35-global-progress}. Thus
the infinite-Force case must serve $C$.

If only finitely many Forces occur, the last version contains
unbounded accepted DATA. Each session has finite quota, so it contains
infinitely many accepted regular sessions. Their strictly increasing
selection heights eventually pass $C$'s permanent-eligibility position
and any earlier cooldown expiry. Cyclic \tagg{NextLive} then cannot pass
$C$ indefinitely: whenever its scan reaches $C$, that region is
eligible. Eligibility changes of other regions cannot make it skip
an eligible $C$ in the fixed cycle. Its selected session must contribute
before contiguous acceptance can pass it. This is again a contradiction.

Consequently every round containing such a debt serves $C$. To see
that a newly permanent-ON $C$ reaches a round, consider an unfinished
earlier round. If its membership never becomes empty, after its finite
removals its nonempty residual set stabilizes. The same infinite-Force
argument forces a further discharge; with only finitely many Forces,
the regular cyclic argument reaches a residual member. Independent
pruning also removes a member rather than prolonging the round. Thus
the old round ends. Unbounded DATA and finite session quotas provide
further duty boundaries; at the next round opening $C$ joins. Repeating
the argument after any one contribution proves repeated service.
\end{proof}

This is accepted regional DATA service. Empty DATA can discharge a
regional debt, so the proposition does not establish transaction
inclusion. Its proof intentionally permits regular sessions outside
the owed set; treating all selections as debt-filtered would prove a
different protocol.

\subsection{Request inclusion and the remaining membership boundary}
\label{sec:v35-request-limit}
The first-DATA REPORT rule repairs the earlier empty-output example
when the regional members are informed. A valid report set has weight
$>F$ and contains a correct member. If every correct member knows a
still-admissible request $x$ absent from the parent, that member reports
$x$ or a request with a smaller finite-predecessor key. A Byzantine
initializer cannot obtain a valid all-EMPTY witness in this situation.
Repeated accepted first blocks must therefore include $x$ or exhaust
that finite predecessor set. \Cref{prop:v38-request-inclusion} gives the
precise conditional argument; one informed signer in each of infinitely
many accepted witnesses is sufficient.

The same-snapshot request relay in \cref{sec:requests} turns one correct
member's receipt into awareness at every correct member of that snapshot.
It remains active after a BFT worker exits. It adds no readiness wait,
changes no public window and makes no whole-pool transfer claim. The
full DATA progress argument above therefore does not depend on a new
request-acquisition service.

This still does not establish request inclusion across every membership
history allowed by the unqualified progress text. Consider disjoint new
rosters for successive slots. A roster can finish its first-DATA reports
before receiving $x$, then receive and relay $x$ after those immutable
reports but before its last service ends. Its old workers and request
owners retain $x$. Its successor is a disjoint roster that can again
report before learning $x$. First-receipt FIFO and a finite local queue
do not prevent this ordering, and a local relay to the old membership
does not deliver to the new one. Repeating this pattern explains why
pointwise eventual arrival is not the awareness property used in the
proof. Whether permanent identity replacement is restricted by the
phrase ``correct leaders receive repeated turns'' also needs its
identity-level meaning; the argument does not silently impose one.

The literal fixed-snapshot communication premise supplies a bound for
old-to-old and new-to-new correct messages. It does not by itself specify
a bound for an old-only member to a new-only member. AUTH authenticates
the roster lineage, but does not transfer pending request bytes. A
physical-region interpretation covering all these endpoints would be
a different clarified communication domain. The present candidate does
not assume that interpretation or add it as a new premise.

\Cref{prop:v39-public-request-import} supplies a local import path
once an original complete public result actually carries the request.
Its continued worker can promote a candidate after ordinary successful
checks against an already prepared exact context, even if the containing
record cannot append. Source visibility alone is not that promotion,
and no import watermark gates Ready or delays the DATA bounds above.
The recurrent informed-attempt trace can leave only one FINAL signature,
so the import rule does not prove formation of the first carrier.

A full repair must implement a request-transfer path under the original
services, or resolve that original model domain and prove the transfer
algorithm against it. A locally retained owner, an unverified snapshot
root, a threshold acknowledgement without all-recipient dissemination,
and a bare-request Espresso submission are not substitutes. No such
handoff wait is inserted into Ready in this version. The environmental
assumptions remain frozen, and the full request-liveness goal remains
open even if the conditional REPORT proof passes review.

\subsection{Implementation and verification scope}
\label{sec:v35-progress-scope}
These arguments concern the stated abstract protocol and the prepared
regional cohort. They require the native adapter to preserve directly
executable anchor state and the \tagg{SOURCESTEP}/ReadyBundle handlers
to implement the existing complete-message processing contract. They
establish neither a general CPU bound nor the same application-validation
deadline for arbitrary observers. Public intrinsic input validity follows
from the honest-FINAL lemma; native output, executable state and
$C_{\rm pub}$ DA correspondence retain their separate sound obligations. They do not certify arbitrary database
imports, archive scans, offline catch-up, or a concrete Espresso profile.
The retry, window, cooldown, health-policy and debt transitions, and
the logical-carrier, public-kernel, byte-retention and per-DATA commit rules, also
require an independent whole-protocol review with the unchanged strict
cutoff and retirement rules. The regional and global progress results
above do not imply the separate awareness condition. The first-DATA
rule has its own conditional request proof, and public import has an
explicit first-carrier premise; arbitrary-roster dissemination
remains unresolved. The logical carrier and original complete-result delivery scope in
\cref{sec:transport-scope} remain essential; no new retry, archive or
authority deadline is inferred from them.
